\section{Mathematical Proofs}
\label{app:theory}

This Online Resource records the formal results used to delimit the empirical
interpretation. The search-theoretic core is established prior art; the exact
posterior-projector presentation remains novelty-unresolved. These proofs have
received a second machine audit, not independent human proof review.

\subsection{Fixed-query phase-sensitive dilution}

Let $\Omega=\{1,\ldots,N\}$ and let $\F$ be uniform over the size-$M$
subsets of $\Omega$, with $\phi=M/N$. Define
\begin{equation}
  \Pi_{\F}=\sum_{x\in\F}|x\rangle\!\langle x|,
  \qquad
  Q_{\F}(\gamma)
  =I+(e^{-i\gamma}-1)(\Pi_{\F}\otimes I_A),
\end{equation}
where $A$ is an arbitrary ancilla register.

\begin{theorem}[Phase-sensitive average bound]
Consider a normalized pure-state query algorithm
\begin{equation}
 |\psi_{\F}\rangle=
 V_qQ_{\F}(\gamma_q)\cdots V_1Q_{\F}(\gamma_1)
 V_0|\psi_{\mathrm{in}}\rangle,
 \label{eq:fixed-query-algorithm}
\end{equation}
where the input, $V_t$, and $\gamma_t$ are independent of the particular
$\F$. If
\begin{equation}
 P_{\F}=\langle\psi_{\F}|(\Pi_{\F}\otimes I_A)|\psi_{\F}\rangle,
\end{equation}
then
\begin{equation}
 \E_{\F}P_{\F}
 \leq
 \min\left\{1,
 \left(1+\sum_{t=1}^{q}|e^{-i\gamma_t}-1|\right)^2\phi
 \right\}.
 \label{eq:phase-sensitive-bound}
\end{equation}
\label{thm:phase-sensitive}
\end{theorem}

\begin{proof}
Replace every query in \cref{eq:fixed-query-algorithm} by the identity and let
$|\varphi_t\rangle$ be the normalized reference state immediately before
slot $t$. It is independent of $\F$. Write
$|\varphi_t\rangle=\sum_x|x\rangle|\alpha_{t,x}\rangle$. Then
\begin{align}
 \E_{\F}\| (\Pi_{\F}\otimes I_A)|\varphi_t\rangle\|^2
 &=\sum_x\Pr[x\in\F]\|\alpha_{t,x}\|^2\\
 &=\frac{M}{N}\sum_x\|\alpha_{t,x}\|^2=\phi.
 \label{eq:reference-overlap}
\end{align}
Let $c_t=|e^{-i\gamma_t}-1|$ and let $d_t(\F)$ be the Euclidean distance
between actual and reference states after slot $t$ and the following
$\F$-independent unitary. Inserting and subtracting the query applied to the
reference state, and using unitarity, gives
\begin{equation}
 d_t(\F)\leq d_{t-1}(\F)
  +c_t\|(\Pi_{\F}\otimes I_A)|\varphi_t\rangle\|.
\end{equation}
The $L^2(\F)$ triangle inequality (Minkowski) and
\cref{eq:reference-overlap} therefore yield
\begin{equation}
  \|d_q\|_{L^2(\F)}\leq\sqrt{\phi}\sum_{t=1}^{q}c_t.
\end{equation}
The final identity-reference state is also independent of $\F$, so its
average marked mass is $\phi$. Projection contractivity gives pointwise
\begin{equation}
 \sqrt{P_{\F}}
 \leq\|(\Pi_{\F}\otimes I_A)|\varphi_{q+1}\rangle\|+d_q(\F).
\end{equation}
A second $L^2$ triangle inequality gives
$\sqrt{\E P_{\F}}\leq(1+\sum_t c_t)\sqrt\phi$. Squaring the nonnegative sides
and applying the probability cap proves \cref{eq:phase-sensitive-bound}.
\end{proof}

The proof is average over random fixed-cardinality subsets. It does not claim
that every particular $\F$ has small overlap with every reference state.
The cases $M=0$ and $M=N$ follow directly, and the cap retains a valid
probability bound at all $q$.

\subsection{Adaptive hard cap, training, and expected queries}

For the phase-flip membership oracle $O_{\F}=I-2\Pi_{\F}$,
$|e^{-i\pi}-1|=2$, so \cref{thm:phase-sensitive} immediately gives the
coarse fixed-schedule form. The following explains the adaptive extension.

\begin{proof}[Proof of the membership-only dilution theorem in the main article]
Introduce registers for the query label, work space, classical transcript,
random seed, discarded environment, halt flag, scratch query, and output.
Replace each measurement/instrument by a Stinespring isometry and retain its
outcome in an orthogonal transcript register. Purify classical randomness and
make feed-forward coherently controlled on the transcript. Mixed initial
states are covered by a purifying environment.

Pad every branch to $q$ query slots. Once a branch halts, coherently swap an
isolated scratch label into the query port, apply the ordinary oracle, and swap
back. The scratch data never recouple to the output. Finally, reversibly copy
the selected classical output label into a designated query register. The
result is a pure $q$-slot algorithm with $\F$-independent interquery unitaries.

For each slot, $\|(O_{\F}-I)|\varphi_t\rangle\|
=2\|(\Pi_{\F}\otimes I)|\varphi_t\rangle\|$. The identity-reference
calculation in \cref{eq:reference-overlap} and the same two Minkowski steps
give
\begin{equation}
 \sqrt{\E P_{\mathrm{success}}}\leq(2q+1)\sqrt\phi.
\end{equation}
Squaring and capping at one proves the theorem. Orthogonal histories also
cover early stopping and feed-forward. A history-controlled direct sum counts
as one query slot only under this explicit controlled membership interface.
\end{proof}

The bit-query oracle
$B_{\F}|x,b\rangle=|x,b\oplus\ind[x\in\F]\rangle$ obeys
$B_{\F}-I=\Pi_{\F}\otimes(X-I)$ and $\|X-I\|=2$, so the same coarse theorem
holds. If slot $t$ instead uses a history-controlled phase $\gamma_t(h)$, the
phase-sensitive result holds with
$c_t=\sup_h|e^{-i\gamma_t(h)}-1|$ under the corresponding counted direct-sum
interface.

If a training procedure satisfies the same membership-only premise, regard
all training circuits, sampled outputs, optimizer transcripts, feedback, and
the final circuit as one adaptive interaction. A deterministic pathwise cap
\begin{equation}
 q_{\mathrm{total}}=q_{\mathrm{training}}+q_{\mathrm{sampling}}
 +q_{\mathrm{feedback}}+q_{\mathrm{final}}
\end{equation}
is substituted for $q$. Classical storage of a query-generated transcript
does not erase its calls. Free explicit coefficients or rich cost values would
instead define a stronger access model.

\subsubsection{Expected-query truncation}
\label{app:expected-query}

Let $Q$ be a nonnegative integer-valued number of membership calls and let
$\bar q=\E Q$, with the expectation covering the subset prior and internal
randomness. For any deterministic integer $T\geq0$, convert executions with
$Q>T$ into failures. The retained protocol has hard cap $T$, and Markov's
inequality on the integer event gives
\begin{equation}
  \Pr(Q>T)=\Pr(Q\geq T+1)\leq\frac{\bar q}{T+1}.
\end{equation}
Therefore
\begin{equation}
 \E P\leq\min\left\{1,
 \frac{\bar q}{T+1}
 +\min\{1,(2T+1)^2\phi\}\right\}.
 \label{eq:expected-query-truncation}
\end{equation}
This is the audited replacement for the invalid substitution
$q=\E Q$ in the main article's hard-cap membership bound. If $Q$ is sampled independently
before the interaction, conditioning additionally yields a bound in terms of
$\E(2Q+1)^2$; that statement is not transferred to transcript-dependent
stopping.

\subsection{Raw-density counterexample and representation padding}

\begin{proof}[Proof of the raw-density counterexample in the main article]
Let
$s=v_0\to v_1\to\cdots\to v_m=t$, assign objective cost and resource one to
every edge, and set budget $m$. The only source--target path uses all $m$
edges and is feasible and optimal. Exactly one of $2^m$ edge selections
encodes it, so $\phistate=2^{-m}$. Starting from $s$, an adjacency-list
algorithm follows the unique outgoing edge at each internal vertex, checks
the additive totals, and writes the $m$ edge identifiers in $O(m)$ time.

If constants $a,c>0$ satisfied both the algorithmic upper bound
$T(I_m)\leq am+a$ and a purported universal lower bound
$T(I_m)\geq c\phistate^{-1/2}=c2^{m/2}$, then
$(am+a)/2^{m/2}\geq c$ for all $m$. The left side tends to zero, a
contradiction for sufficiently large $m$.
\end{proof}

\begin{lemma}[Edge-subdivision representation padding]
Let $I$ be an explicit directed RCSP instance with additive rational
objective coefficients and nonnegative rational resource vectors. Replace one
edge $e=(u,v)$ by a private serial path of $r\geq1$ edges, and assign each
segment $1/r$ of every coefficient of $e$. Then routes are in bijection,
feasibility and objective ordering are preserved, and
\begin{equation}
 \phistate(I^{(r)})=2^{-(r-1)}\phistate(I).
\end{equation}
\label{lem:padding}
\end{lemma}

\begin{proof}
Every new internal vertex has indegree and outdegree one and is incident to no
other edge. A path avoiding $e$ is unchanged. A path using $e$ expands to all
$r$ segments, while any path entering the serial chain must traverse it and
contracts uniquely to $e$. These maps are mutual inverses.

For a route using $e$, the new contribution is
$\sum_{k=1}^{r}c_e/r=c_e$ and likewise
$\sum_{k=1}^{r}a_{e,j}/r=a_{e,j}$ for every resource $j$. Hence paired route
totals, feasibility, objective comparisons, and tied optima are identical.
The number of valid feasible route bitstrings is unchanged, while the raw
edge-bit denominator gains $r-1$ binary variables and hence a factor
$2^{r-1}$.

If a rational coefficient is represented in lowest terms as $p/q$, the
unreduced segment coefficient is $p/(qr)$. Its binary encoding grows by at
most $\lceil\log_2r\rceil+O(1)$. The explicit topology still contains $r$
segments and $r-1$ vertices, so its listing and route-output overhead is
$\Theta(r)$ rather than $O(\log r)$.
\end{proof}

\subsection{Explicit-attribute-query RCSP}

\begin{proof}[Proof of the explicit-attribute-query RCSP theorem in the main article]
Let the known topology contain $K$ two-edge branches
$P_i:s\to u_i\to t$. The first edge of each branch has resource one. The
second has resource one when $z_i=1$ and two when $z_i=0$. With budget two,
\begin{equation}
  P_i\text{ is feasible}\quad\Longleftrightarrow\quad z_i=1.
\end{equation}
Equal fixed objective costs leak no marked index. One branch-feasibility check
uses exactly one query to the attribute oracle defined in the main article, and conversely the
feasibility bit is $z_i$. Finding any feasible route is thus multiple-marked
search over $K$ labels with $M$ marks.

For a uniform size-$M$ array, the membership-only dilution theorem gives
\begin{equation}
 \E_zP_{\mathrm{success}}\leq
 \min\left\{1,(2q+1)^2\frac{M}{K}\right\}.
\end{equation}
Average target success $\tau$ requires
$q\geq\tfrac12(\sqrt{\tau K/M}-1)$. A uniform worst-case guarantee implies
the same average lower bound. In the nontrivial regime $M/K<\tau$, integrality
converts this to $\Omega_\tau(\sqrt{K/M})$. A uniform superposition over
branches plus amplitude amplification uses $O(\sqrt{K/M})$ attribute queries,
giving the matching upper bound.

When $M/K\geq\tau$, a zero-query uniformly random branch already achieves the
target, so the unqualified all-$M$ expression would be false. For exact
success, the order holds for $1\leq M<K$, and $M=K$ needs zero queries. The
input contains $\Theta(K)$ attribute cells and graph entries; no compact
$O(\log K)$ input or RAM-time result is asserted.
\end{proof}

\subsection{Posterior structure and finite classical advice}

Condition on a supported advice value $S=s$ and define
$\overline\Pi_s$, $\lambda_s$, and $\Lambda(S)$ as in the main article.

\begin{lemma}[Posterior projector]
For every supported $s$,
\begin{equation}
 \operatorname{Tr}\overline\Pi_s=M,
 \qquad
 \lambda_s=\max_x\Pr[x\in\F\mid S=s],
 \qquad
 \phi\leq\lambda_s\leq1.
\end{equation}
\label{lem:posterior-projector}
\end{lemma}

\begin{proof}
$\overline\Pi_s$ is diagonal in the label basis and its $x$th diagonal entry
is the posterior inclusion probability. Its trace is the conditional
expectation of $|\F|=M$. The maximum of $N$ nonnegative entries summing to $M$
is at least $M/N$, and every inclusion probability is at most one.
\end{proof}

\begin{proof}[Proof of the posterior structure theorem in the main article]
Fix $s$. Purify measurements, randomness, discarded data, feed-forward, and
early stopping as above. Every identity-oracle reference state
$|\psi_s\rangle$ may depend on $s$ but is independent of the residual
uncertainty in $\F$ conditional on $s$. Consequently,
\begin{align}
 &\E[\langle\psi_s|(\Pi_{\F}\otimes I)|\psi_s\rangle\mid S=s]\\
 &\qquad=\langle\psi_s|(\overline\Pi_s\otimes I)|\psi_s\rangle
 \leq\lambda_s.
 \label{eq:posterior-reference}
\end{align}
Repeating the two conditional $L^2(\F\mid s)$ Minkowski steps with per-slot
phase coefficient $c_t(s)$ gives
\begin{equation}
 \E[P\mid S=s]
 \leq\min\left\{1,
 \left(1+\sum_{t=1}^{q}c_t(s)\right)^2\lambda_s\right\}.
 \label{eq:posterior-conditional}
\end{equation}
Averaging the already capped conditional expression is valid. Since
$c_t(s)\leq2$ and probabilities are at most one,
\begin{align}
 \E P
 &\leq\E_S\!\left[\min\{1,(2q+1)^2\lambda_S\}\right]\\
 &\leq\min\{1,(2q+1)^2\Lambda(S)\}.
\end{align}
The cap is not moved through expectation in the phase-dependent statement;
the displayed hard-cap form is its valid coarser consequence.
\end{proof}

\begin{lemma}[Finite advice alphabet]
If $|\supp S|\leq2^b$, then
\begin{equation}
 \Lambda(S)\leq\min\{1,2^b\phi\}.
\end{equation}
\label{lem:finite-advice}
\end{lemma}

\begin{proof}
For each supported $s$, choose
$x_s\in\arg\max_x\Pr[x\in\F\mid S=s]$. Then
\begin{align}
 \Lambda(S)
 &=\sum_s\Pr(S=s)\Pr(x_s\in\F\mid S=s)\\
 &=\sum_s\Pr(S=s,x_s\in\F)\\
 &\leq\sum_s\Pr(x_s\in\F)
 =|\supp S|\phi\leq2^b\phi.
\end{align}
Duplicate maximizing labels merely repeat a marginal term in an upper bound;
they may loosen, but cannot invalidate, the inequality. The cap by one follows
from $\lambda_s\leq1$.
\end{proof}

Combining the posterior structure theorem in the main article with
\cref{lem:finite-advice} proves the finite-advice inequality stated there. If
target success is $\tau>0$, a necessary integer
condition is
\begin{equation}
 b\geq\max\left\{0,
 \left\lceil\log_2\frac{\tau}{(2q+1)^2\phi}\right\rceil\right\}.
\end{equation}
For $b_{\mathrm{eff}}=\log_2(\Lambda/\phi)$ and
$K_{\mathrm{eff}}=M/\Lambda$,
\begin{equation}
 0\leq b_{\mathrm{eff}}\leq\log_2(1/\phi),\quad
 M\leq K_{\mathrm{eff}}\leq N,\quad
 \frac{N}{K_{\mathrm{eff}}}=2^{b_{\mathrm{eff}}}.
\end{equation}
The target conditions
\begin{equation}
 b_{\mathrm{eff}}+2\log_2(2q+1)\geq\log_2(\tau/\phi),
 \qquad
 K_{\mathrm{eff}}\leq\frac{(2q+1)^2M}{\tau}
\end{equation}
are necessary only. A candidate set of size $K$ gives
$\lambda_s=M/K$ only when the conditional posterior is uniform over size-$M$
subsets of that candidate set.

\subsection{One-shot finite-dimensional quantum advice}

% CLAIM: C-T10
For completeness, the audited extension covers one pre-search advice state,
not a refreshed or interactive structure provider.

\begin{proposition}[One-shot quantum-advice dimension bound]
Suppose one $\F$-dependent density operator $\rho_{\F}$ on a total accessible
advice register of Hilbert-space dimension $d$ is supplied before an otherwise
membership-only fixed-query interaction. Then
\begin{equation}
 \E P_{\mathrm{success}}
 \leq\min\left\{1,
 \left(1+\sum_t c_t\right)^2d\phi\right\},
\end{equation}
and hence the coarse hard-cap factor is $(2q+1)^2d\phi$.
\label{prop:quantum-advice}
\end{proposition}

\begin{proof}
For an $\F$-independent trace-preserving reference channel $A$ acting before a
query, $\rho_{\F}\preceq I_d$ implies
\begin{align}
 \E_{\F}\operatorname{Tr}[(\Pi_{\F}\otimes I)A(\rho_{\F})]
 &\leq\E_{\F}\operatorname{Tr}[(\Pi_{\F}\otimes I)A(I_d)]\\
 &=\phi\operatorname{Tr}A(I_d)=d\phi.
\end{align}
A canonical purification of each $\rho_{\F}$ preserves the marked projected
norm; an inaccessible purifying reference does not enlarge the counted advice
dimension. The phase-sensitive hybrid proof then replaces the reference
overlap $\phi$ by $d\phi$. Multiple simultaneously supplied advice registers
count by their joint accessible dimension.
\end{proof}

Mixed advice and entanglement with an inaccessible purifier are included.
Fresh copies, repeatedly refreshed advice, and an interactive
$\F$-dependent provider are outside \cref{prop:quantum-advice}.

The assumptions and corresponding interpretation boundaries are consolidated
in \cref{tab:proof-assumptions}.

\begin{table}[ht]
\centering
\caption{Assumptions that bound the interpretation of the formal results.}
\label{tab:proof-assumptions}
\small
\begin{tabularx}{\textwidth}{@{}lX@{}}
\toprule
Result & Required assumptions \\
\midrule
Fixed/adaptive dilution & Uniform fixed-cardinality prior; all remaining feasible-set dependence enters through counted membership queries; pathwise hard cap for adaptive protocols \\
Trained specialization & Training, sampling, feedback, and final membership calls are counted jointly \\
Expected-query form & Integer truncation with an explicit tail term; no substitution of $\mathbb{E}Q$ for a hard cap \\
Posterior structure & One classical pre-search channel; conditional reference circuit independent of residual feasible-set uncertainty \\
Classical advice & Advice support has cardinality at most $2^b$; $b$ is not mutual information \\
Quantum advice & One accessible pre-search state of total dimension $d$; no refreshing or interactive provider \\
Attribute-query RCSP & Explicit length-$K$ attribute array accessed only through counted quantum random access \\
\bottomrule
\end{tabularx}
\end{table}

